DensePose tells a downstream model \emph{where} a pixel lies on the body, not
\emph{how the surface is turned} there. We add the missing rotational degrees of
freedom: a DINOv3+DPT network predicting, per pixel and in one pass, the usual part
index and chart coordinates \emph{and} a full $SO(3)$ surface frame --- to our
knowledge the first orientation output for DensePose. Anchored to the body template
rather than the chart, the frame is continuous across chart seams, and composing it
forward transports any template-attached direction into the image: camera-space
normals, hair and weave tangents, and the twist about the normal that a normal map
cannot carry. The model reaches $0.741$ geodesic-point-similarity (GPS) AP on
DensePose-COCO minival under a ground-truth-box protocol, above every model in our own
lineage --- though we re-score no published system under this harness, and the oriented
output costs $2.1$ AP points at matched training step before the recipe recovers them.
Retraining a CVPR'25 pose-guided generator with our correspondence model as the one
changed input improves paired LPIPS by $0.0116$, better on $68.7\%$ of $8{,}512$ pairs
and of the same order as the largest published step on that benchmark; FID does not
move. Predicting a frame correctly is also what forced us to repair the correspondence:
the official DP24 charts are non-injective on roughly half the hand and foot surface at
a median $3$D collision of $60$--$67$\,mm, which a sheet split removes; we supervise
correspondence in metres through a differentiable chart-to-template operator ($-46\%$
hand and foot surface error against a matched control); and we add occlusion data under
a sequence-level holdout ($14.80^\circ\!\rightarrow\!6.89^\circ$ frame error there).
